\section{Discarded techniques and decisions}\label{sec:tecnicas-descartadas}

Besides its final result, the factor analysis (\cref{sec:analisis-factorial})
produced three methodological choices that were tried and explicitly
rejected, and two variables that the model could not accommodate. Documenting
them prevents them from looking like omissions: each was discarded for a
concrete, verifiable reason, not for convenience.

\subsection{Horn's parallel analysis}

This is the standard criterion for deciding how many factors to retain, and
here it retains seven ---more than twice the number given by Velicer's MAP,
which was finally adopted. It is not followed, because the result is an
artifact of sample size rather than a substantive signal: the simulated
threshold lies at 0.06 on the reduced matrix and at 1.06 on the full matrix.
With the 16,028 training observations of the 14-predictor block, the
eigenvalues that pure noise would produce concentrate very tightly around
their expected value, the simulation band narrows almost to a point, and any
barely positive eigenvalue exceeds it. The argument for discarding it holds by
looking at the simulated thresholds alone, without examining a single factor
loading. Velicer's MAP is used instead because it is a descriptive minimum
that does not include $n$ in its formula and does not inflate with sample
size.

\subsection{Maximum likelihood extraction}

The extraction of \cref{sec:analisis-factorial} uses iterated principal axes
rather than maximum likelihood, which requires multivariate normality. The
Box-Cox transformations of \cref{sec:transformaciones} brought the marginals
close to symmetry ($|\text{skewness}|<0.5$) but did not make them normal, so
that assumption does not hold for these data. The principal axis solution was
validated against the reference implementation of the R package
\texttt{psych} on the same block: it matches up to the order and sign of the
factors, with communalities equal to the third decimal.

\subsection{Oblique rotation}

Before adopting varimax, an oblique rotation (oblimin), which allows the
factors to correlate, was tried. It found correlations of up to 0.32 between
them ---physically reasonable, since wind disperses pollutants and is not
independent of the combustion load. It is rejected for the instrumental use of
the scores: the goal of \cref{sec:clasificacion_tecs} is to have
non-collinear predictors, and an oblique solution reintroduces exactly the
multicollinearity problem that factor analysis was meant to solve. The
correlation between factors is reported as a finding but is not carried into
the classification model.

\subsection{Two variables the factor model could not accommodate}

\texttt{NOX} and \texttt{RH\_min\_diurna} produced communalities greater than
1 ---the so-called \textbf{Heywood cases}, impossible because communality is a
proportion of variance. \texttt{NOX} is, by chemical definition, the sum of
\texttt{NO} and \texttt{NO2} ($r=0.985$ against that sum): it is not redundant
in an ordinary statistical sense but an accounting identity that the PCA of
\cref{sec:pca} absorbed without complaint and that factor analysis cannot
split into communality and uniqueness. \texttt{RH\_min\_diurna} correlates at
0.892 with \texttt{RH\_media}; unlike \texttt{NOX} it is not an identity, so
\cref{sec:analisis-factorial} kept it under observation until the second fit,
with the full diagnosis, forced the same exit. Both were removed from the
predictor block, which is left with 13 variables and a maximum VIF of 2.54.
That PCA flagged neither problem is itself a result: the factor model exposes
redundancies that a total-variance decomposition has no way of reporting.
